\section{Threats to Validity}

\paragraph{Construct validity.}
Security and functionality are defined by independently frozen evaluators.
Code validity, Oracle support, unknown coverage, oracle-evaluable secure-code
yield, functionality, and joint success are therefore reported separately.
The resulting claims concern these measured outcomes rather than unrestricted
program security.

\paragraph{Internal validity.}
The Prompt TSG defines what can be edited; observational discovery determines
what should be tested; blocked randomization identifies the assigned
prompt-policy effect. Typed constraints and FCI organize the experimental
budget but do not convert Prompt TSG or PAG relations into effect estimates.
Prospective hypothesis fixation, arm-complete eligibility before assignment,
blocked randomization, complete assigned-arm accounting, and independent
measurement protect the randomized contrast. Post-assignment fidelity and
non-target-drift diagnostics cannot redefine that contrast.

\paragraph{Selection and multiplicity.}
The shared candidate universe, top-$K$ budget, expected directions, outcomes,
contrasts, and multiplicity families are fixed before held-out outcomes.
Selector failures remain in the $K$ denominator, while simultaneous inference
resamples the highest-level semantic task clusters. This design keeps the
selector comparison and the hypothesis-specific effect claims within their
prospectively declared evidence families.

\paragraph{External validity.}
Each effect remains scoped to its context query, editable factor, operation,
task and CWE population, prompt realization distribution, generator model,
and evaluators. Cross-realization or cross-model claims require their separate
frozen support rules. RQ4 likewise concerns expert-perceived utility under its
sample and presentation design; it does not establish objective repair
accuracy or validate the computational effects.
